\subsection{Propósito e Arquitetura}
\label{sec:proposito_arquitetura}

O manipulador robótico realiza tarefas de captura, manipulação e acoplamento durante as operações de RVD/B.
Após a aproximação da espaçonave visitante, o MR atua na fase final de acoplamento, realizando movimentos de alinhamento e fixação entre o efetuador $\{E\}$ e a porta de acoplamento do alvo.
Nessa etapa, o controle orbital é substituído por uma dinâmica de contato, em que o manipulador e a base passam a interagir mecanicamente por meio das forças de reação internas.

Por estar em um ambiente de microgravidade, a conservação do momento angular total do sistema acopla o movimento da base e o do manipulador, portanto, cada movimento das juntas do MR gera torques de reação que afetam a atitude da base, exigindo que o controle de orientação compense essas alterações.
A compensação é implementada por estratégia de controle acoplado, as quais a cinemática e a dinâmica do MR são incorporadas às leis de controle de atitude.

A configuração adotada para o MR é do tipo 4R1P, sendo quatro juntas revolutas que realizam o posicionamento e orientação da ferramenta, e uma junta prismática, que fornece extensão linear ao efetuador.
Além disso, o projeto do MR leva em conta as restrições físicas e operacionais, como os limites de curso das juntas, capacidade máxima de torque e a velocidade dos atuadores.
